\documentclass[conference]{IEEEtran}
\IEEEoverridecommandlockouts
\usepackage{cite}
\usepackage{amsmath,amssymb,amsfonts}
\usepackage{algorithmic}
\usepackage{graphicx}
\usepackage{textcomp}
\usepackage{xcolor}
\usepackage{booktabs}

\def\BibTeX{{\rm B\kern-.05em{\sc i\kern-.025em b}\kern-.08em%
    T\kern-.1667em\lower.7ex\hbox{E}\kern-.125emX}}

\begin{document}

\title{Assignment 1: Classifiers}

\author{\IEEEauthorblockN{Kevin Deshayes}
\IEEEauthorblockA{\textit{DV2638} \\
\textit{Blekinge Institute of Technology}\\
Karlskrona, Sweden \\
kede23@student.bth.se}
}

\maketitle

\section{Introduction}
The Wine Quality dataset contains physicochemical features
describing physicochemical properties of wine and a
discrete quality label. This work focuses on predicting wine quality
using supervised learning for the red-wine subset (1599 samples), where
quality labels range from 3 to 8.

\section{Method}
An 80/20 stratified split was used to preserve class distribution.
Features were standardized using the \texttt{StandardScaler}. Two
classifiers were trained:

\begin{itemize}
    \item \textbf{Logistic Regression}: lbfgs solver, max\_iter=1000.
    \item \textbf{Random Forest}: 1000 trees, random\_state=2000.
\end{itemize}

Repeated 3-Fold Cross-Validation (10 repeats) was applied using
macro-F1 due to imbalance. The best model was retrained on the full
training split and evaluated on the held-out test set. Oversampling was
also tested to analyze imbalance effects.


\section{Results}
\subsection{Cross-Validation}
Random Forest achieved higher macro-F1 and lower variance.

\begin{table}[htbp]
\centering
\caption{Cross-Validation Performance}
\begin{tabular}{lcc}
\toprule
\textbf{Model} & \textbf{Mean F1} & \textbf{Std} \\
\midrule
Logistic Regression & 0.289 & 0.029 \\
Random Forest & \textbf{0.320} & \textbf{0.010} \\
\bottomrule
\end{tabular}
\end{table}

\subsection{Test Performance}
\begin{table}[htbp]
\centering
\caption{Test Set Results}
\begin{tabular}{lcc}
\toprule
\textbf{Model} & \textbf{Accuracy} & \textbf{Macro F1} \\
\midrule
Random Forest (unbalanced) & \textbf{0.678} & \textbf{0.401} \\
Random Forest (balanced) & 0.622 & 0.363 \\
\bottomrule
\end{tabular}
\end{table}

Balancing improved minority-class recall during training but reduced test
generalization since the real test distribution remained imbalanced.

\section{Conclusion}
Random Forest showed the strongest validation and test performance.
Oversampling improved validation metrics but decreased test generalization,
thus the unbalanced model is preferred. Dataset imbalance remains the main
limiting factor.


\end{document}
